% A.69 candidate: not integrated pending independent proof/source audit.
\subsection{The nonstable dihedral family at every two-power level}

\begin{statement}{A.69}{Proved}{Every two-power member of the nonstable dihedral family}
Let \(m\ge2\) be a power of two, and let \(C_m\) be the smooth proper
normalization of
\[
                    v^2=w^{2m}-6w^m+1.
\]
Give \(C_m\setminus S_m\), where \(S_m\) consists of the four points
above \(w=0,\infty\), the actual degree-\(2m\) map to
\(E\setminus\{O\}\) below. Its proper Jacobian has potential good
reduction of \(2\)-rank zero at two. Every cusp difference is killed
by \(m\). Consequently both all-prime assertions of Task A.5 hold
for this cover and for every actual intermediate of its dihedral
Galois closure: every induced cusp packet has precisely its selected
inertia as the fixed geometric pro-\(\ell\) subgroup under every
open derived Galois subgroup, and every arithmetically normalized
procyclic pro-\(\ell\) subgroup is cuspidal. The same assertions hold
for the plus-six family, identified with this one after a finite
extension by \(w\mapsto\zeta_{2m}w\).

This proves a uniform family and all its intermediates; the universal
cover question A.45 remains open.
\end{statement}

\begin{proof}
All curves, automorphisms, quotient maps, and chosen cusp sections can
be defined over one finite number-field extension. Further finite
extensions used for reduction are harmless in the formulation of
A.41--A.42. We always distinguish the dihedral deck action from the
additional cyclic automorphism used to control the Jacobian.
At the chosen place over two we enlarge this field so that the
elliptic base has good reduction, the Jacobians below have good
reduction, and all displayed endomorphisms and automorphisms are
defined. The local actions throughout fix the resulting coefficient
valuation ring.

\paragraph{The actual cover.}
The equation is smooth of genus \(m-1\), and the actual map is
\[
 X=\frac{(w^m-1)^2}{4w^m},\qquad
 Y=\frac{v(w^{2m}-1)}{8w^{3m/2}},\qquad Y^2=X^3-X.
                                                        \tag{A.69.1}
\]
It factors through elliptic duplication after
\[
 x=\frac{w^m-1+v}{2},\qquad y=w^{m/2}(x-1),
 \qquad y^2=x^3-x,\qquad w^{m/2}=\frac{y}{x-1}.
                                                        \tag{A.69.2}
\]
The divisor of the last elliptic function is
\(T_0+T_{-1}-T_1-O\). Hence the first map is cyclic of degree
\(m/2\), totally ramified at those four points, and unramified
elsewhere; duplication is unramified. Thus (A.69.1) has degree
\(2m\), its four cusp indices are \(m/2\), and it is finite
\'{e}tale off \(S_m\). Its actual deck transformations
\[
 r(w,v)=(\zeta_m w,-v),\qquad s(w,v)=(1/w,-v/w^m)
                                                        \tag{A.69.3}
\]
satisfy \(r^m=s^2=1\), \(srs=r^{-1}\), and give \(2m\)
distinct transformations. The deck group is therefore
\(D_{2m}\). For \(m=2\) this is the unramified degree-four
elliptic isogeny case, so \(J_{C_2}\) is isogenous to \(E\).

\paragraph{Potential good reduction: the precise ramified input.}
After a finite extension at two the original elliptic curve has a
smooth proper model. The branch divisor consists of the single
section \(O\), so its closure is \'{e}tale over the valuation ring.
The finite dihedral group is nilpotent. The ramified
nilpotent-cover theorem, as stated precisely in
Lehr--Matignon, Theorem 6.13, applies: the branch closure is \'{e}tale
over the valuation ring, and generic tameness means tameness in
characteristic zero, which holds here. It gives potential good
reduction of \(J_{C_m}\), equivalently a tree as the dual graph of
the proper stable fibre. Sa\"idi's statement of the same ramified
extension, printed page 662, confirms this scope.
\cite{fghLehrMatignon2005,fghSaidi2007}
We use only this tree/potential-good conclusion. Raynaud's
unramified theorem about strictly interior slopes of contracted
components is not being extended to a marked branch fibre.

\paragraph{The new factor and its integral characteristic-two restriction.}
Put
\[
 \mu(w,v)=(\zeta_mw,v),\qquad
 \sigma=\mu^{m/2},\qquad C_m/\langle\sigma\rangle=C_{m/2}
                                                        \tag{A.69.4}
\]
for \(m\ge4\); the quotient coordinate is \(w^2\).
Let \(A_m\) be the anti-invariant isogeny factor of \(J_{C_m}\)
for \(\sigma\). Then
\[
 J_{C_m}\sim J_{C_{m/2}}\times A_m,\qquad
 \dim A_m=m/2,
 \qquad \mathbf Q(\zeta_m)\hookrightarrow\operatorname{End}^0(A_m).
                                                        \tag{A.69.5}
\]
Indeed the regular differentials \(w^{j-1}dw/v\),
\(1\le j\le m-1\), have \(\mu\)-eigenvalues \(\zeta_m^j\).
The anti-invariant indices are exactly the odd indices; since \(m\)
is a power of two, every such eigenvalue is primitive of order
\(m\). This proves the field action, including its unit, and shows
that \(\mu^k-1\) is invertible on every rational realization of
\(A_m\) whenever \(m\nmid k\).

After the indicated finite extension giving good reduction, all
these endomorphisms extend to the abelian scheme and the same
unital field acts on the rational \'{e}tale Tate space of the special
\(2\)-divisible group. Since
\(\mathbf Q(\zeta_m)\otimes\mathbf Q_2\) is a field of degree
\(\varphi(m)=m/2\), the special \(2\)-rank of \(A_m\) is either
zero or \(m/2\). In the latter case \(A_m\) is ordinary.
This is a dimension argument in the actual special \'{e}tale Tate
space, not a conclusion from a Hodge character pattern.

We prove by induction that the ordinary alternative cannot occur.
The base \(m=2\) was identified with an isogeny of the original
supersingular elliptic curve. Assume the result for \(m/2\), and
suppose for a contradiction that \(A_m\) is ordinary.

\paragraph{Componentwise consequences of an ordinary anti-invariant factor.}
Choose the stable marked model of \((C_m,S_m)\) over a complete
valuation ring \(R\), after the finite extension needed to define
all automorphisms and sections. Its residue field may be replaced
by an algebraic closure. Its dual graph \(\Gamma\) is a tree;
inserting marked rational components does not change this fact or
the product decomposition of its identity-component Picard scheme.
Thus the good special Jacobian is the product of the Jacobians of
the positive-genus components. The involution \(\sigma\) acts on
this actual product.

If \(\sigma\) exchanges two positive-genus components, their
diagonal and anti-diagonal factors are both isogenous to the same
component Jacobian. The diagonal is a factor of the specialization
of \(J_{C_{m/2}}\), whose \(2\)-rank is zero, while the
anti-diagonal is a factor of the ordinary \(A_m\). This is
impossible. A component on which \(\sigma\) acts identically
contributes only to the invariant factor; no restriction on its
geometric genus is inferred from that trivial action.

Consider a component \(B\) fixed by \(\sigma\), on which it acts
nontrivially. Its quotient \(B'=B/\langle\sigma\rangle\) is a
smooth proper curve, and its invariant Jacobian is isogenous to
\(J_{B'}\), by the actual pullback and norm identities. Hence
\(f(B')=0\), where \(f\) denotes \(2\)-rank. The anti-invariant
Prym is a factor of \(A_m\) and is ordinary whenever it has
positive dimension.

The Deuring--Shafarevich formula and Riemann--Hurwitz give, if
\(r_B\) is the number of fixed points and \(d_P\) the local
different exponent,
\[
 f(B)=r_B-1,\qquad
 g(B)-g(B')=g(B')-1+\frac12\sum_{P\ {
m fixed}}d_P,
 \qquad d_P\ge2.                                  \tag{A.69.6}
\]
There must be a fixed point, since an unramified connected
double cover of a \(2\)-rank-zero curve would give negative
\(2\)-rank. If the Prym had dimension zero, (A.69.6) would force
\(g(B')=0\), \(r_B=1\), and \(g(B)=0\); it contributes no
positive genus. Otherwise ordinarity says
\(f(B)=g(B)-g(B')\). Consequently
\[
      r_B=g(B')+\frac12\sum d_P
      \quad\Longrightarrow\quad
      g(B')=0,\qquad d_P=2\ \hbox{for every fixed }P.
                                                        \tag{A.69.7}
\]
Thus every positive-genus component on which \(\sigma\) is
nontrivial is ordinary, and \(\sigma\) is its hyperelliptic
involution. The Deuring--Shafarevich input is the ordinary
curve formula, not a formula applied to the marked vertical fibre.
\cite{fghCrew1984,fghYang2024}

\paragraph{An ordinary hyperelliptic curve has no cyclic fourth root of its involution.}
Here is the required elementary ramification argument. If a
\(2\)-group acts faithfully on an ordinary smooth curve, apply
Deuring--Shafarevich and Riemann--Hurwitz and subtract twice the
first formula from the second. The result is
\[
 0=2|G|\bigl(g(B/G)-f(B/G)\bigr)
       +\sum_{P\in B/G}\frac{|G|}{e_P}
                          \bigl(d_P-2(e_P-1)\bigr).
                                                        \tag{A.69.8}
\]
Every term is nonnegative: for wild \(2\)-group inertia the
zeroth and first ramification groups coincide, so
\(d_P=\sum_{i\ge0}(|G_i|-1)\ge2(e_P-1)\).
Indeed the local minimal polynomial of a uniformizer has derivative
\(\prod_{g\ne1}(t-g(t))\); its valuation is
\(\sum_{g\ne1}\operatorname{ord}_t(t-g(t))\), which is this
ramification-group sum. We use the classical lower numbering,
as in \cite[\S2.2, equation (3) and Exercise 2.12(c)]{fghBouwWewers}.
Equality forces the second ramification
group to be trivial. The first ramification quotient embeds in
the additive residue field via the coefficient in
\(t\mapsto t+a t^2+\cdots\); it has exponent two. Thus no
cyclic inertia group of order four can fix a point of an ordinary
curve.

Suppose an automorphism \(a\) of an ordinary curve has order four
and \(a^2=\sigma\), with \(B/\sigma\simeq\mathbf P^1\).
Its induced involution on \(\mathbf P^1\) has a unique fixed
point \(Q\) in characteristic two. If \(Q\) were unramified in
the double cover, \(a\) would permute its two points, so \(a^2\)
would fix them individually; but the deck involution \(\sigma\)
exchanges them. Thus \(Q\) is ramified, its unique point above it
is fixed by \(a\), and the cyclic inertia of order four
contradicts (A.69.8).

\paragraph{The primitive cyclic action forces one full orbit of elliptic components.}
Take a positive-genus component \(B\) contributing to the
anti-invariant factor, and let \(d\) be its orbit length under
\(\mu\). The restriction of \(\mu^d\) to \(B\) has exact order
\(m/d\). Otherwise a smaller power of \(\mu\) would be the
identity on the nonzero Jacobian contribution of this entire
orbit, contrary to the primitive field action in (A.69.5).
Moreover \(\sigma\) acts nontrivially on \(B\), so \(m/d\ge2\).
The induced automorphism on \(B/\sigma=\mathbf P^1\) has order
\(m/(2d)\). A finite \(2\)-power-order element of
\(\operatorname{PGL}_2\) in characteristic two has order at most
two. Hence \(m/d\) is two or four; the fourth-root argument
excludes four. Therefore \(d=m/2\).

The sum of the genera of these components is precisely
\(\dim A_m=m/2\): exchanged positive-genus pairs were excluded,
and each invariant quotient here has genus zero. There is
therefore exactly one orbit of \(m/2\) components, each an
ordinary elliptic curve. On each of them \(\sigma\) has exactly
two fixed points by (A.69.6).

\paragraph{A smooth fixed point really lifts in this wild situation.}
Suppose \(P\) is a smooth special point fixed by \(\sigma\) on
a component on which \(\sigma\) is nontrivial. The action fixes
the coefficient ring \(R\). In a smooth formal coordinate its
completed local ring is \(R[[t]]\), and
\(F(t)=\sigma(t)-t\) has nonzero reduction of finite order.
Weierstrass preparation writes \(F\) as a unit times a
distinguished monic polynomial. Thus the fixed locus in this
formal disc is finite flat of positive degree over \(R\), and
has generic fixed points specializing to \(P\). The smooth
formal-disc identification realizes its roots as points on the
generic curve, and properness extends them to sections after a
finite extension. This is the usual
one-variable Weierstrass preparation theorem, with no uniformizer
factor because the reduction is nonzero.
\cite[Corollary 2.16]{fghVonk2021}
For the ordinary
elliptic components the degree is two, since their different
exponent is two. There is no vertical factor: nonzero reduction
of \(F\) is exactly what rules it out.

The generic fixed points of \(\sigma(w,v)=(-w,v)\) are exactly
the four cusps. Each of those cusps is fixed individually by
\(\mu\). Therefore a smooth \(\sigma\)-fixed point cannot lie
on one of the nontrivially permuted elliptic components: its
lifted cusp section would force that component to be fixed by
\(\mu\). This argument concerns smooth special points only;
no lifting assertion is made at a node or on a component on
which \(\sigma\) is the identity.

\paragraph{The rooted-tree genus contradiction.}
Root \(\Gamma\) at the component containing one selected marked
cusp. This vertex is \(\mu\)-fixed. All the \(m/2\) ordinary
elliptic vertices have the same positive distance from this root.
On each such vertex \(B\), the parent edge and its attaching
point are fixed by \(\sigma\), because \(\sigma\) fixes both
the root and \(B\). This accounts for one of its two fixed
points. Its other fixed point cannot be smooth by the preceding
paragraph, so it is a different node, lying on a child edge.

The descendant subtree beyond that child edge contains no marked
cusp: if it did, the path from the root to that \(\mu\)-fixed
mark would be fixed pointwise as a graph by \(\mu\), forcing
\(B\) itself to be fixed. It contains no member of the ordinary
elliptic orbit, since every member of that orbit has the same
distance from the root. A nonempty finite unmarked subtree
attached by one edge must contain a positive-genus component;
otherwise it has a rational leaf with fewer than three special
points, contrary to stability of the marked model. These
descendant subtrees are disjoint for the distinct ordinary
elliptic vertices, again because those vertices have equal
distance from the root. Their contributions, together with
the elliptic vertices themselves, give
\[
             g(C_m)\ \ge\ \frac m2+\frac m2=m,
                                                        \tag{A.69.9}
\]
contrary to \(g(C_m)=m-1\). This excludes the ordinary
alternative and completes the induction proving special
\(2\)-rank zero.

\paragraph{Cusps and every actual intermediate.}
The identities
\[
 \begin{aligned}
 (v-w^m+1)(v+w^m-1)&=-4w^m,\\
 (v-w^m-1)(v+w^m+1)&=-8w^m
 \end{aligned}                                           \tag{A.69.10}
\]
have exactly the cusp divisors used in A.68 with eight replaced
by \(m\). More explicitly, write \(P_{0,\pm}\) for the points
with \(v=\pm1\), and \(P_{\infty,\pm}\) for those with
\(v/w^m\to\pm1\). The expansion
\(\sqrt{1-6u^m+u^{2m}}=1-3u^m-4u^{2m}+O(u^{3m})\)
at zero, and at infinity with \(u=w^{-1}\), gives
\[
 \begin{aligned}
 \operatorname{div}(v-w^m+1)&=mP_{0,-}-mP_{\infty,-},\\
 \operatorname{div}(v+w^m-1)&=mP_{0,+}-mP_{\infty,+},\\
 \operatorname{div}(v-w^m-1)&=mP_{0,+}-mP_{\infty,-},\\
 \operatorname{div}(v+w^m+1)&=mP_{0,-}-mP_{\infty,+}.
 \end{aligned}                                             \tag{A.69.12}
\]
These connect all four cusps and show that every cusp difference
is annihilated by \(m\).
The cusp inertia in (A.69.3) is
\(I=\langle r^2\rangle=\Phi(D_{2m})\), including the trivial
Frattini subgroup when \(m=2\). It is normal. For every proper
subgroup \(H<D_{2m}\), the actual intermediate
\(D=C_m/H\) has
\[
        \#S_D=\#(H\backslash D_{2m}/I)
                 =[D_{2m}:HI]\ge2.                 \tag{A.69.11}
\]
For completeness, if \(HI=D_{2m}\), surjectivity onto
\(D_{2m}/I\simeq C_2^2\) supplies in \(H\) an odd rotation
and a reflection. The odd rotation generates \(\langle r\rangle\),
so these two elements generate \(D_{2m}\) and \(H=D_{2m}\).
This proves the displayed strict inequality without an assumption
about chosen lifts. The actual compact curve \(D\) maps finitely to
\(E\), so has genus at least one. Its Jacobian is an isogeny
quotient of \(J_{C_m}\) by the actual norm, hence has potential
good reduction and special \(2\)-rank zero. Over an algebraic
closure the pushforward of a point is its image with coefficient
one; ramification indices do not enter this pushforward.
Taking norms of the functions in (A.69.10) therefore shows that
every difference of cusps on \(D\) is killed by \(m\).

Their reductions in the good special Jacobian are zero, because
a \(2\)-rank-zero abelian variety has no nonzero geometric
\(2\)-power torsion points. The compact-type Abel-path criterion
of A.41 gives a rational joining path to a distinct companion
for every cusp. A.42 therefore proves the fixed-subgroup
assertion at every coefficient prime, including two. The
normalizer assertion follows from A.59, since every such
closure is a \(2\)-group. The case \(H=D_{2m}\) is the already
proved original punctured elliptic curve. This completes the
proof for every intermediate and both signs.
\end{proof}
